\documentclass[11pt,a4paper]{amsart}
\usepackage[T1]{fontenc}
\usepackage{lmodern}
\usepackage[margin=25mm]{geometry}
\usepackage{amsmath,amssymb,amsthm,microtype}
\usepackage[hidelinks]{hyperref}
\hypersetup{
  pdftitle={Erdos Problem 928 as a corollary of OpenAI's joint Dickman law},
  pdfauthor={Jiho Bae},
  pdfsubject={An explicit implication of OpenAI's joint Dickman law}
}
\newtheorem{corollary}{Corollary}
\newcommand{\N}{\mathbb N}
\title{Erd\H{o}s Problem 928 as a corollary of\\
OpenAI's joint Dickman law}
\author{Jiho Bae}
\date{October 9, 2026}

\begin{document}
\begin{abstract}
We identify Erd\H{o}s Problem~928 as a consequence of OpenAI's joint
Dickman law. An elementary counting estimate shows that the different
inequality and threshold conventions do not change the natural density.
\end{abstract}
\maketitle

Let $P(n)$ be the largest prime divisor of the integer $n\ge2$, and let
$\rho$ denote the Dickman function. For fixed $\alpha,\beta\in(0,1)$
and real $X\ge2$, put
\begin{align*}
 A_{\alpha,\beta}(X)
   &=\#\{n\in\N:2\le n\le X,\ P(n)<n^\alpha,
                         \ P(n+1)<(n+1)^\beta\},\\
 B_{\alpha,\beta}(X)
   &=\#\{n\in\N:2\le n\le X,\ P(n)\le n^\alpha,
                         \ P(n+1)\le n^\beta\}.
\end{align*}
Problem~928 in Bloom's collection~\cite{Bloom928} asks whether
$A_{\alpha,\beta}(X)/X$ has a limit, and records the predicted value
$\rho(1/\alpha)\rho(1/\beta)$.

OpenAI's joint Dickman law~\cite[Theorem~1.1, p.~2]{OpenAI2026}
(result family~012 of its mathematics release) gives
\begin{equation}\label{eq:JD}
 \lim_{X\to\infty}\frac{B_{\alpha,\beta}(X)}X
   =\rho(1/\alpha)\rho(1/\beta).
\end{equation}
\begin{corollary}\label{cor:928}
The joint Dickman law \eqref{eq:JD} implies that, for every
$\alpha,\beta\in(0,1)$,
\[
  \lim_{X\to\infty}\frac{A_{\alpha,\beta}(X)}X
   =\rho(1/\alpha)\rho(1/\beta).
\]
\end{corollary}

\begin{proof}
An integer counted by $B_{\alpha,\beta}(X)$ but not by
$A_{\alpha,\beta}(X)$ must satisfy $P(n)=n^\alpha$, since
$n^\beta<(n+1)^\beta$. On this exceptional set, the map $n\mapsto P(n)$
is injective and takes positive integer values at most $X^\alpha$.
There are therefore at most $\lfloor X^\alpha\rfloor$ such integers.

Conversely, an integer counted by $A_{\alpha,\beta}(X)$ but not by
$B_{\alpha,\beta}(X)$ must satisfy
\[
 n^\beta<P(n+1)<(n+1)^\beta.
\]
The open intervals $(n^\beta,(n+1)^\beta)$ are pairwise disjoint as
$n$ varies. Thus $n\mapsto P(n+1)$ is injective on this exceptional
set, and its values are positive integers below $(X+1)^\beta$.
Its cardinality is at most $\lfloor(X+1)^\beta\rfloor$.
The two bounds give, for every $X\ge2$,
\begin{equation}\label{eq:error}
 |A_{\alpha,\beta}(X)-B_{\alpha,\beta}(X)|
   \le \lfloor X^\alpha\rfloor+\lfloor(X+1)^\beta\rfloor.
\end{equation}
Dividing by $X$, the right-hand
side tends to zero because $\alpha,\beta<1$; the conclusion follows
from~\eqref{eq:JD}.
\end{proof}

The higher-order independence conjecture for three or more consecutive
integers mentioned in~\cite{Bloom928} is outside the scope of this corollary.

\begin{thebibliography}{9}
\bibitem{Bloom928}
T.~F. Bloom, \emph{Erd\H{o}s Problem \#928}, Erd\H{o}s Problems,
\url{https://www.erdosproblems.com/928}, accessed October~9, 2026.

\bibitem{OpenAI2026}
OpenAI, \emph{The joint Dickman law for consecutive integers},
preprint, September~24, 2026, result family~012; Theorem~1.1, p.~2.
\href{https://github.com/openai/math/blob/fd4aeeb2ee4fc729c18d98444fed42fd0529eeeb/preprints/The-joint-Dickman-law-for-consecutive-integers-September-24-2026/paper.pdf}{OpenAI Math Release},
repository version \texttt{fd4aeeb2ee4f}.
\end{thebibliography}

\end{document}
